% Industry resume, Grover Aruquipa
% Compile with pdfLaTeX (also works on Overleaf). One column on purpose:
% applicant tracking systems read single-column text far more reliably than sidebars.
\documentclass[10pt,a4paper]{article}
\usepackage[T1]{fontenc}
\usepackage[utf8]{inputenc}
\usepackage[scaled=0.95]{helvet}
\renewcommand{\familydefault}{\sfdefault}
\usepackage[top=0.9cm,bottom=0.7cm,left=1.4cm,right=1.4cm]{geometry}
\usepackage{enumitem}
\usepackage{titlesec}
\usepackage{xcolor}
\usepackage{microtype}
\usepackage[hidelinks]{hyperref}

\definecolor{ink}{HTML}{111111}
\definecolor{accent}{HTML}{0D2B61}
\definecolor{muted}{HTML}{555555}
\color{ink}
\pagestyle{empty}
\setlength{\parindent}{0pt}
\setlength{\parskip}{0pt}

\titleformat{\section}{\normalsize\bfseries\color{accent}}{}{0em}{\MakeUppercase}[\vspace{-6pt}{\color{accent}\rule{\linewidth}{0.6pt}}]
\titlespacing*{\section}{0pt}{5pt}{2pt}
\setlist[itemize]{leftmargin=1.1em,itemsep=0.5pt,topsep=1pt,parsep=0pt,label=\textbullet}

% \role{Title}{Organization}{Location}{Dates}
\newcommand{\role}[4]{\textbf{#1}, #2 \hfill {\color{muted}#3 \textbar\ #4}\par}

\begin{document}

%================ HEADER ================
{\LARGE\bfseries Grover Aruquipa}\par\vspace{2pt}
{\large Robotic Controls Engineer \textbar\ Real-Time Control, Robot Bring-Up, Sim-to-Real}\par\vspace{3pt}
{\small Qu\'ebec, QC, Canada \ \textbullet\ \
\href{mailto:grover.aruquipa.9@gmail.com}{grover.aruquipa.9@gmail.com} \ \textbullet\ \
\href{https://www.linkedin.com/in/grover-aruquipa-996917126}{linkedin.com/in/grover-aruquipa-996917126} \ \textbullet\ \
\href{https://groveraruquipa.github.io}{groveraruquipa.github.io}}\par\vspace{1pt}
{\small\bfseries Available from January 2027 \ \textbullet\ \ Ready to relocate to Z\"urich and work on site}\par
% TODO: add a current phone number on the line above if you want it on the CV.

%================ SUMMARY ================
\section{Summary}
Robotics engineer finishing a PhD in Mechanical Engineering (Universit\'e Laval, defense early 2027). I have brought up
three manipulators with integrated grippers (4, 7 and 8 actuators), from the kinematic and dynamic model through
fabrication and wiring to 500\,Hz real-time control and debugging on the bench. I write the URDF and MJCF models of my robots,
validate them and every trajectory in MuJoCo before hardware, and I am training a reinforcement learning policy in NVIDIA
Isaac Sim for a gripper robot designed to sense grasp force through its motor currents. Previously worked in Z\"urich at
Hitachi ABB Power Grids R\&D (machine learning on sensor data).

%================ SKILLS ================
\section{Technical Skills}
\begin{itemize}
\item \textbf{Control and estimation:} multi-axis real-time control at 500\,Hz (OPAL-RT, Simulink Real-Time), trajectory
generation and tracking, controller tuning; kinematic and dynamic modelling, Jacobians, redundancy resolution, singularity
analysis; force estimation from motor currents (in development).
\item \textbf{Simulation and learning:} MuJoCo, NVIDIA Isaac Sim; robot description files written by me (URDF, MJCF);
reinforcement learning (Isaac Sim, Q-learning); PyTorch, TensorFlow; OpenCV.
\item \textbf{Software:} Python, C++ (real-time control), object-oriented programming, MATLAB and Simulink; ROS and ROS~2,
MoveIt, Pinocchio; Linux, Git; logging and analysis scripts for robot experiments.
\item \textbf{Robot hardware and bring-up:} manipulators and grippers, electric motors, drives and encoders, first power-up
and commissioning, root-cause debugging across mechanics, electronics and software; SolidWorks, Altium.
\end{itemize}

%================ EXPERIENCE ================
\section{Experience}
\role{Doctoral Researcher}{Robotics Laboratory, Universit\'e Laval}{Qu\'ebec, Canada}{2023 -- present}
\begin{itemize}
\item Designed, built and controlled three parallel manipulators (4, 7 and 8 actuators) with every motor
mounted on the base and the gripper driven through the mechanism, enabling unlimited rotation of the grasped object.
\item Implemented real-time multi-motor control at 500\,Hz; validated a 4-DoF coaxial wrist-gripper (0.25\,kg moving mass,
180\,mm footprint) over ten continuous tool revolutions with sub-degree tracking.
\item Derived closed-form inverse kinematics, Jacobians and singularity conditions, and developed a redundancy
strategy that keeps unlimited rotation free of collisions and singularities.
\item Training a reinforcement learning grasping policy in NVIDIA Isaac Sim for a Delta robot with an integrated Sarrus
gripper that I am building, aimed at sensing the grasp force through the motor currents; deployment on the real robot is
the next step.
\item Write the URDF and MJCF description of each robot and validate it, with every trajectory, in MuJoCo before running on hardware.
\item Lead two graduate students day to day on an 8-actuator (6+2)-DoF robot; joint paper under review at ASME J.\
Mechanisms and Robotics.
\item First author of 3 journal papers (ASME J.\ Mechanisms and Robotics, Mechanism and Machine Theory) and 5 more
under review; reviewer for ASME J.\ Mechanisms and Robotics.
\end{itemize}

\vspace{3pt}
\role{Visiting Researcher}{Institute of Robotics, Johannes Kepler University Linz}{Linz, Austria}{Apr. -- May 2026}
\begin{itemize}
\item Developed grasp-quality indices for manipulators with integrated grasping (Mitacs Globalink Research Award).
\end{itemize}

\vspace{3pt}
\role{Research Intern}{FEMTO-ST (CNRS)}{Besan\c{c}on, France}{2022 -- 2023}
\begin{itemize}
\item Modeled a redundant dexterous parallel robot with two unlimited rotations for micro-manipulation.
\item Wrote real-time control software in C++ for the robot experiments.
\end{itemize}

\vspace{3pt}
\textbf{Industry roles in vision and machine learning} \hfill {\color{muted}2020 -- 2021}\par
\begin{itemize}
\item \textbf{Computer Vision Engineer}, Sencon Srl (remote, 2021): computer vision for industrial anomaly detection.
\item \textbf{Graduate Research Intern}, Hitachi ABB Power Grids R\&D (Zurich, 2020): machine learning on sensor data from high-voltage circuit breakers.
\item \textbf{Publication:} reinforcement learning (Q-learning) for algorithmic trading, IEEE ESCI 2021.
\item \textbf{Data Scientist}, KTH Royal Institute of Technology (remote, 2021): data analysis for environmental impact assessment.
\end{itemize}

%================ EDUCATION ================
\section{Education}
\textbf{PhD, Mechanical Engineering}, Universit\'e Laval \hfill {\color{muted}Qu\'ebec, Canada \textbar\ 2023 -- early 2027 (expected)}\par
{\small Kinematically redundant parallel robots with unlimited rotations and integrated grasping. Advisors: C.~Gosselin, P.~Lambert.}\par\vspace{2pt}
\textbf{MSc, Smart Integrated Systems}, Universit\'e Bourgogne--Franche-Comt\'e \hfill {\color{muted}Besan\c{c}on, France \textbar\ 2021 -- 2023}\par\vspace{2pt}
\textbf{BEng, Mechatronics Engineering}, Bolivian Catholic University \hfill {\color{muted}La Paz, Bolivia \textbar\ 2014 -- 2019}\par
{\small Thesis: reconfigurable Delta robot for product picking.}

%================ AWARDS ================
\section{Awards and Languages}
\begin{itemize}
\item 1st Place, Graduate Division, ASME Student Mechanisms and Robotics Design Competition (2025), team lead.
\item Finalist, MassRobotics competition, Boston, with one of my robot prototypes, team lead.
\item Mitacs Globalink Research Award (2026); 1st Place, Student Paper Competition, Argentine Robotics Conference (2019).
\item Languages: Spanish (native), English (full professional), French (intermediate).
\end{itemize}


%================ PUBLICATIONS
\section{Selected Publications}
{\small First author on all papers. Full list at groveraruquipa.github.io.}\par\vspace{2pt}
\begin{itemize}
\item \textbf{G.~Aruquipa}, P.~Lambert, C.~Gosselin, ``A compact coaxial parallel wrist-gripper with unlimited tool-axis rotation,'' under review, \emph{Robotics and Autonomous Systems}.
\item \textbf{G.~Aruquipa}, P.~Lambert, C.~Gosselin, ``Analysis and design of a novel (6+1)-DoF parallel robot with unlimited rotation and grasping capabilities,'' \emph{Mechanism and Machine Theory}, vol.~224, art.~106438, 2026.
\item \textbf{G.~Aruquipa}, V.~Muralidharan, P.~Lambert, C.~Gosselin, ``Analysis and design of a novel compact four-degree-of-freedom parallel robot with unlimited rotation,'' \emph{ASME Journal of Mechanisms and Robotics}, vol.~18, no.~6, art.~061001, 2026.
\item \textbf{G.~Aruquipa}, P.~Lambert, R.~Dahmouche, C.~Gosselin, ``Parallel robots for grasping tasks: a review,'' \emph{ASME Journal of Mechanisms and Robotics}, vol.~18, no.~12, art.~120801, 2026. doi:10.1115/1.4072612
\item \textbf{G.~Aruquipa}, P.~Lambert, C.~Gosselin, ``Kinematic analysis and design of a novel 9-DoF parallel robot with grasping capabilities,'' \emph{CCToMM Symposium on Mechanisms, Machines and Mechatronics}, 2025, pp.~117--127.
\end{itemize}

%================ REFERENCES ================
\section{References}
\begin{itemize}
\item \textbf{Prof.~Cl\'ement Gosselin}, Full Professor, Department of Mechanical Engineering, Universit\'e Laval, Qu\'ebec,
Canada. PhD supervisor. \href{mailto:clement.gosselin@gmc.ulaval.ca}{clement.gosselin@gmc.ulaval.ca}
\item \textbf{Prof.~Patrice Lambert}, Assistant Professor, Department of Mechanical Engineering, Universit\'e Laval,
Qu\'ebec, Canada. PhD co-supervisor. \href{mailto:patrice.lambert@gmc.ulaval.ca}{patrice.lambert@gmc.ulaval.ca}
\item \textbf{Univ.-Prof.~Dr.-Ing.~Andreas M\"uller}, Head of the Institute of Robotics, Johannes Kepler University Linz,
Austria. Host during my research stay (2026). \href{mailto:a.mueller@jku.at}{a.mueller@jku.at}
\end{itemize}

\end{document}
